%!TEX root = ../dokumentation.tex

\section{Einleitung}

Moderne Sprachmodelle wie ChatGPT bestechen durch eine verblüffend menschliche Konversationsfähigkeit. Fragt man diese Systeme jedoch direkt nach ihrer Vernunft, geben sie oft selbst an, lediglich statistische Textwahrscheinlichkeiten zu berechnen und Intelligenz nur zu simulieren \citeM{i}{WeilChat2025}{1}{1}. Daraus ergibt sich eine zentrale philosophische und praktische Forschungsfrage: Verbirgt sich hinter dieser technischen Simulation in der Praxis nicht doch eine funktionale Form echter diskursiver Vernunft?

Um diese Frage zu beantworten, reicht der klassische Turing-Test, welcher lediglich auf erfolgreiche Täuschung und Imitation abzielt nicht aus. Diese Hausarbeit unterzieht das Modell ChatGPT stattdessen einem \enquote{verschärften Turing-Test}. Ziel ist es zu prüfen, ob die KI fähig ist, Entscheidungen normativ zu begründen, Argumente qualitativ zu bewerten, eigene Positionen rational zu korrigieren und ihr Verhalten auf einer Meta-Ebene zu reflektieren, also insgesamt in der Lage ist, Vernunft aufzuweisen.

